\section{Conclusiones}
\subsection{Cinemática directa}
\pendiente{Se sintetizará la relación entre las coordenadas articulares y la pose del efector, a partir del modelo desarrollado y sus comprobaciones.}
\subsection{Cinemática inversa}
\pendiente{Se resumirán las soluciones obtenidas, las condiciones de alcanzabilidad y los errores de reconstrucción.}
\subsection{Cinemática diferencial}
\pendiente{Se expondrá la correspondencia entre velocidades articulares y cartesianas, sustentada en la comparación del Jacobiano con derivadas numéricas.}
\subsection{Singularidades}
\pendiente{Se identificarán las configuraciones singulares y sus consecuencias para el movimiento del efector.}
\subsection{Control cinemático}
\pendiente{Se evaluará el seguimiento de trayectoria mediante errores, velocidades y comportamiento cerca de singularidades.}
\subsection{Conclusión general}
\pendiente{Se integrarán los hallazgos de los cinco temas y se precisarán el alcance y las limitaciones de los resultados obtenidos.}
